\documentclass[11pt]{article}
\usepackage[margin=1in]{geometry}
\usepackage{amsmath,amsthm,amssymb,mathtools}
\usepackage[hidelinks]{hyperref}
\newtheorem{theorem}{Theorem}
\newtheorem{lemma}[theorem]{Lemma}
\newcommand{\E}{\mathbb E}
\newcommand{\PP}{\mathbb P}
\newcommand{\Sg}{\mathfrak S}
\newcommand{\ind}{\mathbf 1}
\title{Empty and crowded residue classes at the first-prime scale}
\author{Meiji \quad | \quad Erd\H os tournament, round 119}
\date{24 November 2026}
\begin{document}
\maketitle
\begin{abstract}
For every fixed positive constant $C$, a positive proportion of the
reduced classes modulo every sufficiently large integer $Q$ contain
no prime below $C\phi(Q)\log Q$. For every fixed $c>0$ and integer $m$,
a positive proportion contain at least $m$ primes below $c\phi(Q)\log Q$.
We prove both conclusions by transporting the signed divisor-square
method for ordinary short intervals to an array of rows with growing
common difference $Q$. The method uses only one detected prime at a
time, not a prime-tuple conjecture. We neither determine a limiting
distribution nor assert that either positive proportion is uniform in
its parameters.
\end{abstract}

\section{Statement and prior work}
Write $N_Q(a,x)=\#\{p\le x:p\text{ prime},\ p\equiv a\pmod Q\}$ and
$p(a,Q)=\min\{p:p\equiv a\pmod Q\}$. All proportions below give equal
weight to reduced residues. The constants may depend on the fixed
parameters but never on the factorization of $Q$.
\begin{theorem}\label{main}
For every fixed $C>0$, $c>0$ and positive integer $m$ there exist
$d_C,e_{c,m}>0$ such that for all sufficiently large $Q$,
\begin{align*}
 \#\{a\in(\mathbb Z/Q)^\times:p(a,Q)> C\phi(Q)\log Q\}&\ge d_C\phi(Q),\\
 \#\{a\in(\mathbb Z/Q)^\times:N_Q(a,c\phi(Q)\log Q)\ge m\}&\ge e_{c,m}\phi(Q).
\end{align*}
\end{theorem}
Erd\H os \cite[p.~57]{E49} predicts failure of an almost-all upper
bound for every constant multiple, proving a late-class result for
one multiple $1+c_0$ along infinitely many moduli and the early result
for $m=1$. His 1965 formulation \cite[p.~209, (33)]{E65}, now numbered
971, asks for one $1+c_0$ over all large moduli. The latter statement
already has public proofs claimed by Han (July 2026) and Li (September
2026) \cite{Han,Li}; we claim no priority for the catalogue statement.
Leung \cite{Leung} obtains a Poisson law under a uniform prime-tuple
hypothesis. A result concerning the maximum over classes, such as
\cite{LiPr}, does not give a positive proportion. Conversely, the first part gives
\[
 \lim_{Q\to\infty}\frac{\max_{(a,Q)=1}p(a,Q)}
                             {\phi(Q)\log Q}=+\infty:
\]
for every fixed $C$ there is at least one late class for all large $Q$.
This is the weaker all-modulus conjecture separated by Granville and
Pomerance \cite[p.~194]{GP} from their stronger $\phi(Q)\log^2 Q$
lower bound, which we do not prove. Primorial moduli are included.

The analytic model is OpenAI's signed multidimensional sieve for
ordinary prime gaps \cite{OAI}. We use its arithmetic-free
cancellation mechanism with attribution. The changing step here is
essential: in the physical array its prime divisors affect every row
simultaneously. We remove them before averaging the singular series,
and show that the remaining normalization cancels with the number of
rows, even for very smooth $Q$. This note contains the transfer and
the two applications, without requiring the ordinary-gap theorem.

\section{A mixed-moment transfer}
Fix $\lambda>0$, let $L=\log Q$, $u=Q/\phi(Q)$,
$h=\lfloor\lambda L/u\rfloor$, and $I=\{1,\ldots,h\}$. We use the
standard bound $u\ll\log\log Q$, hence $h\to\infty$ and $hu/L\to\lambda$.
Throughout $a$ is uniform on the representatives $1\le a\le Q$ with
$(a,Q)=1$. Put $n_b=a+bQ$ and $\vartheta(n)=\ind_{n\;\mathrm{prime}}\log n$.
For fixed smooth tests $F_r$ in dimensions $j_r$, take compactly
supported smooth extensions from the nonnegative orthants to the full
spaces. On the orthant $F_r$ vanishes for $\sum x_i>\rho_r$. Define
\[
 D_{F_r}(a;b_1,\ldots,b_{j_r})=
 \sum_{d_i\mid n_{b_i}}\prod_i\mu(d_i)\,
 F_r((\log d_i/L)_i).
\]
A zero-dimensional test is a scalar. Across all tests prescribe a
partition of the coordinates $\Gamma$ into $t$ nonempty fibers
$M_1,\dots,M_t$, meaning coordinates in a fiber have the same row;
different fibers use distinct rows $b_1,\dots,b_t\in I$.
If $\delta=1$, allow one additional prime-marked row $b_0\in I$
distinct from these; if $\delta=0$ there is no mark. Write $J$ for
the set of $t+\delta$ row indices, $Y_0=1$, $Y_1=\vartheta(n_{b_0})$,
and
\[
 \Sg_Q(J)=\prod_{p\nmid Q}
 (1-\nu_p(J)/p)(1-1/p)^{-|J|},\qquad
 \nu_p(J)=\#(J\bmod p).
\]
\begin{lemma}\label{mixed}
Suppose $\rho=\sum_r\rho_r<1$ if $\delta=0$, and $\rho<1/2$ if
$\delta=1$. For fixed dimensions, tests and fiber partition, uniformly
over the distinct row indices,
\[
 \E\,Y_\delta\prod_r D_{F_r}
 =u^{t+\delta}L^{-t}
 \big\{\Sg_Q(J)\mathcal C+O(L^{-1/2}(\log L)^A)\big\}.
\]
Here $A$ is fixed and $\mathcal C$ depends only on the tests and
the partition. In particular no relative error is asserted when a local
obstruction makes $\Sg_Q(J)=0$.
\end{lemma}
\begin{proof}
Expand the divisor sums. All surviving tuples are squarefree in each
coordinate, have $\prod_{\gamma\in\Gamma}d_\gamma\le Q^\rho$,
and are coprime to $Q$. Write $D=\mathrm{lcm}(d_\gamma)$.
For each $p\mid D$, all its coordinates must have row index congruent
to one common class $b\bmod p$. This imposes one residue on $a$ modulo
$D$. Without a mark, its conditional density is
$D^{-1}+O(2^{\omega(Q)}/\phi(Q))$: expand $\ind_{(a,Q)=1}$ by
M\"obius inversion and apply CRT for each divisor of $Q$. The number of
tuples is $O(Q^\rho(1+L)^{|\Gamma|})$, by successive harmonic summation.
Since $2^{\omega(Q)}=Q^{o(1)}$, the total error is $O(Q^{\rho-1+o(1)})$.

With a mark, if a selected row $b$ agrees with $b_0\bmod p$, then
$p\mid n_{b_0}$ and the marked term is zero, since $n_{b_0}>Q>D$.
Otherwise $n_{b_0}$ is in a reduced class $r\bmod D$ within
$(b_0Q,(b_0+1)Q]$. A prime in this interval already ensures $(a,Q)=1$.
Thus the conditional marked mass is $u/\phi(D)$ plus an error.
For clarity, write $E_D$ for the maximum, over these intervals and
all reduced residues modulo $D$, of the absolute difference between the
sum of $\vartheta$ and $Q/\phi(D)$. For each fixed $K\ge1,B>0$,
\[
 \sum_{D\le Q^\rho}\mu^2(D)K^{\omega(D)}E_D\ll QL^{-B}.
\]
Indeed the maximal-endpoint Bombieri--Vinogradov theorem up to
$X=(h+1)Q\ll QL$ gives the unweighted estimate with any log saving,
since $Q^\rho\le X^{1/2}(\log X)^{-D_0}$ for every fixed $D_0$ eventually.
Take differences of two endpoint estimates. Replacing von Mangoldt
weights by prime-only weights costs at most
$O(\sqrt X(\log X)^2)$ per modulus, hence $O(Q^{1/2+\rho+o(1)})$ even
with weights. The trivial count in an interval gives
$E_D\ll Q L^2/D$ ($D<Q^\rho<Q$, $D/\phi(D)\ll\log L$).
Therefore $\sum_{D\le Q^\rho}\mu^2(D)K^{2\omega(D)}E_D
\ll Q L^{O_K(1)}$, by the Euler product for the harmonic sum.
Cauchy--Schwarz between this bound and the unweighted estimate proves
the display. At most $2^{|\Gamma|\omega(D)}$ tuples have lcm $D$;
bounded tests and division by $\phi(Q)$ absorb their total error to
arbitrary negative powers of $L$. Only the moduli $D$, not $QD$,
enter this application of Bombieri--Vinogradov.

Write $w_\gamma=1+is_\gamma$, $z_\gamma=w_\gamma/L$. Smooth extension
and Fourier inversion give
$F_r(\mathbf x)=\int\Phi_r(\mathbf s)\exp(-\sum w_\gamma x_\gamma)d\mathbf s$
on the orthant, with Schwartz $\Phi_r$ (the Fourier normalization
is included in $\Phi_r$). The main local-density sum, before Fourier
inversion, was finite by support; one may equivalently extend it to
all tuples since the tests vanish outside that support.
After inversion the sum converges absolutely, uniformly in $\mathbf s$,
by $\prod_p(1+O_{|\Gamma|}(p^{-1-1/L}))\ll L^{O_{|\Gamma|}(1)}$.
There is an external factor $u^\delta$ and the local factors are
$P_p=1$ for $p\mid Q$ and otherwise
\[
 P_p=1+\frac1{p-\delta}
 \sum_{\varnothing\ne U\subseteq\Gamma}^{*}
 (-1)^{|U|}p^{-z_U},\quad z_U=\sum_{\gamma\in U}z_\gamma.
\]
The star requires all rows in $U$ to agree modulo $p$ and, if marked,
to avoid the marked row modulo $p$.

Separate the universal polar product via $H_p=P_p/B_p$, where
\[
 B_p=\prod_{i=1}^t\prod_{\varnothing\ne U\subseteq M_i}
 (1-p^{-1-z_U})^{-(-1)^{|U|}}.
\]
For $|s_\gamma|\le\sqrt L$, the pole expansion of $\zeta(1+z_U)$
gives
\[
 \prod_pB_p=L^{-t}\mathcal K(\mathbf s)(1+O(L^{-1/2})),\qquad
 \mathcal K=\prod_{i=1}^t\prod_{\varnothing\ne U\subseteq M_i}
 \left(\sum_{\gamma\in U}w_\gamma\right)^{(-1)^{|U|+1}}.
\]
Indeed the sum of $(-1)^{|U|}$ over nonempty subsets of one fiber is
$-1$. The kernel grows at most polynomially; the complementary Fourier
tail costs $O(L^{-B})$ for any fixed $B$ by the absolute Euler bound.

At $\mathbf z=0$, $B_p=(1-1/p)^t$. For each distinct allowed row
class the subset sum is $-1$; with a mark the number of allowed classes
is $\nu_p(J)-1$. It follows that
\[
 \prod_p H_p(0)=u^t\Sg_Q(J).
\]
We justify evaluation at zero even at locally obstructed tuples.
For $p\le L^2$, the factors $H_p(z)$ and $H_p(0)$ have modulus at most
$1+C/p$ (increase the fixed $C$ for small primes), and their
difference is $O(L^{-1/2}\log p/p)$ on the clipped Fourier region.
This follows by differentiating their finite local formulas along
the straight segment from $0$ to $z$; all denominators
$1-p^{-1-z_U}$ are bounded away from zero there. Telescoping the products
costs $O(L^{-1/2}(\log L)^A)$ by Mertens' estimates. For $p>L^2$
larger than every row difference, the first-order terms in $P_p$ and
$B_p$ agree if $p\nmid Q$, so $H_p=1+O(p^{-2})$. For $p\mid Q$ we only have
$1+O(1/p)$, but $\sum_{p\mid Q,\,p>L^2}1/p\le L/(L^2\log(L^2))$.
Both tails are negligible, also when multiplying the small-prime
products, which are at most $(\log L)^{O(1)}$. Thus the difference
from the zero-value product is bounded additively as above.
Integrating leaves
$\mathcal C=\int (\prod_r\Phi_r)\mathcal K\,d\mathbf s$.
Since $u\ge1$, normalization gives the asserted error.
\end{proof}

\begin{lemma}\label{av}
For every fixed $l$, over ordered distinct $l$-tuples of indices in $I$
the mean of $\Sg_Q(J)$ is $1+o(1)$ as $Q\to\infty$.
\end{lemma}
\begin{proof}
The empty and singleton cases are exact. Put $y=(\log h)/4$ and
$P=\prod_{p\le y,p\nmid Q}p\le h^{1/2}$ by the prime number theorem
for all large $h$. Over independent uniform residues modulo $p$,
the mean of $1-\nu_p/p$ is $(1-1/p)^l$ by introducing one independent
test residue and asking it to avoid all $l$ coordinates. CRT gives
mean one for the truncated product over complete residue boxes mod $P$.
Reduction of $I$ mod $P$ changes total variation by $O_l(P/h)$;
repeats have probability $O_l(1/h)$ and the product is bounded
by $(\log y)^{O_l(1)}$. For distinct indices let
$\Delta=\prod_{i<j}|b_i-b_j|\le h^{\binom l2}$.
The logarithms of tail factors are $O_l(1/p^2)$ off $\Delta$,
$O_l(1/p)$ on it, for $p>y>2l$.
Their sum is $O_l(1/y+\log\Delta/(y\log y))=o(1)$ uniformly,
whether or not primes dividing $Q$ are omitted. Thus the tail is
$1+o(1)$ multiplicatively, which also preserves any zero.
\end{proof}
\section{Transfer to a quadratic form}
Let $\tau=1/8$. Fix a finite family: a real scalar $f_0$ and real
symmetric smooth $f_j$ compactly supported in
$\{t_i>0,\sum t_i<\tau\}$ for $1\le j\le k$. Set $f_{k+1}=0$ and
\[
 F_j(\mathbf x)=\int_{t_i\ge x_i}f_j(\mathbf t)d\mathbf t,\quad
 (Tf_{j+1})(\mathbf t)=\int_0^\infty f_{j+1}(\mathbf t,s)\,ds,
 \quad F_0=f_0.
\]
The cumulative tests extend smoothly and compactly (integrate after
zero extension and multiply by a smooth cutoff on the negative side).
On the orthant $(-1)^j\partial_1\cdots\partial_jF_j=f_j$, zero on
every coordinate face for $j>0$. Define
\[
 Z=\sum_{j=0}^k\sum_{S\subset I,\,|S|=j}D_{F_j}(a;S),\qquad
 V=L^{-1}\sum_{b\in I}\vartheta(n_b),\qquad \alpha_j=\lambda^j/j!.
\]
Write $N=\#\{b\in I:n_b\text{ prime}\}$; since
$Q<n_b\le(h+1)Q$, we have $N\le V\le(1+O(\log L/L))N$.

\begin{lemma}\label{forms}
For each fixed family as $Q\to\infty$,
\begin{align*}
 \E Z^2&\to w:=\sum_{j=0}^k\alpha_j\|f_j\|_2^2,&
 \E Z^4&=O(1),\\
 \E Z^2 V&\to v:=\lambda\sum_{j=0}^k\alpha_j
       \|f_j+Tf_{j+1}\|_2^2,& \E V^4&=O(1).
\end{align*}
\end{lemma}
\begin{proof}
In a product of two subset tests every fiber has size one or two.
Its factor in $\mathcal K$ for size one is $w_1$; for size two it is
$w_1w_2/(w_1+w_2)=\int_0^\infty w_1w_2 e^{-(w_1+w_2)t}dt$.
Fourier inversion therefore expresses $\mathcal C$ as the integral
of the two signed complete derivatives: a shared coordinate gets
a common integration variable, and any unpaired coordinate is set to zero.
This follows by absolute integration, since the Fourier transforms are
Schwartz and real parts of all $w_i$ equal one. Thus off-diagonal terms
have $\mathcal C=0$, and equal subsets of size $j$ have
$\mathcal C=\int |f_j|^2$.

For each fixed coincidence pattern with $t$ distinct rows there are
$O(h^t)$ labelings and Lemma~\ref{mixed} gives total unmarked error
$O((hu/L)^t L^{-1/2}(\log L)^A)=o(1)$. The same holds for a distinct
marked row after division by $L$, using $O(h^{t+1})$ labelings.
A label constraint from ordering subset coordinates can be dropped
when bounding nonnegative singular-series terms; for diagonal subsets
the average is precisely the symmetric average in Lemma~\ref{av}.
Hence the diagonal terms contribute
$\binom hj(u/L)^j(1+o(1))\|f_j\|_2^2$.
This proves the first identity. For four factors the kernel integrals
are finite fixed constants. By the same counting and nonnegative
singular-series average each pattern contributes $O(1)$, proving the
fourth moment bound ($4\tau<1$).

To mark a row $b_0$, observe that when $n_{b_0}$ is prime its only
allowed divisor in a supported test is $1$: the alternative exceeds
$Q$. Split every subset according to membership of $b_0$. On the
marked event, $Z$ thus becomes the sum over subsets of
$I\setminus\{b_0\}$ with the new test
$\widetilde F_j=F_j+F_{j+1}(\cdot,0)$ in dimension $j$, including
dimension zero. Its signed complete derivative is $f_j+Tf_{j+1}$,
which still vanishes on the faces in positive dimensions.
Lemma~\ref{mixed} applies with $2\tau<1/2$ and a distinct mark.
There are $h\binom{h-1}{j}$ diagonal terms. Averaging over the mark
as well, the counting and normalization now tend to
$\lambda^{j+1}/j!$. This gives the claimed $v$.

Finally take a smooth compact univariate test $G$ with $G(0)=1$
and $G(x)=0$ on $x\ge1/16$. Pointwise
$N\le\sum_{b\in I}D_G(a;b)^2$: for primes the sum is exactly one
in that row and every summand is nonnegative. The fourth power on the
right expands into eight tests whose combined budget is at most $1/2$.
There are finitely many row coincidence types; each of their kernel
constants is bounded and each singular-series average is bounded
as above. Thus $\E N^4=O(1)$, and comparability gives $\E V^4=O(1)$.
\end{proof}

\section{Forcing the two tails}
For completeness we give the arithmetic-free suppression calculation of
\cite[Prop.~4.2]{OAI} and its one-level amplification counterpart.
For any $M>0$ choose a nonnegative $g\in C_c^\infty((0,\infty))$
such that $\int g^2=1$, $\int g=\sqrt M$, and
$\mu_g:=\int t g(t)^2dt<\beta$, with fixed $0<\beta<\gamma<\tau$.
Indeed take a smooth cutoff of $1/t$ equal to $1/t$ on $[1,R]$,
bounded by $1/t$ and supported in $[1/2,2R]$. Normalize its $L^2$
norm and dilate (preserving that norm) to prescribe its integral.
The new mean is at most $M\log(4R)/(\log R)^2$, small for large $R$.
Let $B$ exceed the support endpoint of $g$ and choose smooth $0\le\chi\le1$
equal to one up to $\beta$ on the orthant, zero at and above $\gamma$.
Set
\[
 P_j=k^{j/2}\prod_{i=1}^j g(kt_i),\qquad S_j=\sum_{i=1}^jt_i .
\]
Then $\|P_j\|_2=1$ and $TP_{j+1}=\sqrt{M/k}P_j$.
Under the density $P_j^2$, $S_j$ has mean $j\mu_g/k\le\mu_g$
and variance at most $\operatorname{Var}_{g^2}(t)/k$ for $j\le k$.
Consequently $\PP_j(S_j>\beta-B/k)=O_g(1/k)$ uniformly. Outside this
exception all cutoffs are one, even throughout a deletion integral
over the full support of $g$.

\emph{Suppression.} Set $M=\lambda$ and $r=\lfloor\sqrt k\rfloor$.
On the $r$ levels $k-r+1\le j\le k$ take
$f_j=(-1)^j(r\alpha_j)^{-1/2}P_j\chi(S_j)$, and zero on the rest.
These are interior-supported test functions. Then
$w=1-O_g(1/k)$. On an interior level the expression
$\sqrt{r\alpha_j}(f_j+Tf_{j+1})/((-1)^jP_j)$, interpreted on the
support of $P_j$, is
\[
 \chi(S_j)-\sqrt{(j+1)/(\lambda k)}
          \int g(y)\chi(S_j+y/k)\,dy .
\]
Its absolute value is at most one everywhere and at most $r/k$ off
the exceptional event. At each of the two ends the normalized
contribution to $v/\lambda$ is at most $1/r$ (at the lower end use
$\alpha_{j-1}/\alpha_j=j/\lambda$ and nonnegativity).
Thus $v/\lambda\le 2/r+(r/k)^2+O_g(1/k)\to0$. We fix one such family
with $w>0$ and $v<w/2$. Since $\ind_{N\ge1}\le V$,
\[
 \liminf_{Q\to\infty}\E[Z^2\ind_{N=0}]\ge w-v>0 .
\]
By Cauchy--Schwarz and the bound for $\E Z^4$, a fixed positive
fraction of reduced classes has $N=0$ for all sufficiently large $Q$.

\emph{Amplification.} Keep only
$f_k=\alpha_k^{-1/2}P_k\chi(S_k)$, with $M>0$ to be chosen.
Then $w\to1$ as $k\to\infty$. With $P_{k-1}^2$-probability
$1-O_g(1/k)$, deletion gives exactly
$Tf_k=\alpha_k^{-1/2}\sqrt{M/k}\,P_{k-1}$; everywhere its absolute
value is at most that expression. Hence
\[
 \lambda\alpha_{k-1}\|Tf_k\|_2^2\longrightarrow M,
 \qquad v/w\longrightarrow \lambda+M.
\]
Given $m$, fix $M$ and then $k$ so that $v>(m+1)w>0$. Only afterwards
let $Q\to\infty$. On $N<m$, $V\le m(1+o(1))$, so
$\liminf\E[Z^2V\ind_{N\ge m}]\ge v-mw>0$.
H\"older's inequality bounds that expectation above by
$(\E Z^4)^{1/2}(\E V^4)^{1/4}\PP(N\ge m)^{1/4}$.
Lemma~\ref{forms} gives a fixed positive lower bound for this probability.

For the first conclusion of Theorem~\ref{main}, take $\lambda>C$.
Then $hQ\ge C\phi(Q)L$ eventually. All candidate primes up to this
height in a given class have index $0\le b\le h$; among the positive
proportion of empty rows discard the $a$ which are themselves prime.
There are at most $\pi(Q)\ll Q/L=o(\phi(Q))$ of these. For the second,
take $\lambda=c/2$. Since $Q/(\phi(Q)L)=u/L\to0$, each detected prime
is at most $(h+1)Q<c\phi(Q)L$ eventually. This completes both proofs.

\paragraph{Boundary of the result.}
The two assertions concern different collections of classes.
They give neither an asymptotic frequency for any fixed count nor
a tail bound as $C$ or $m$ increases with $Q$. No generalized
Riemann hypothesis or prime-tuple assumption is used. The catalogue problem
971 has weaker earlier public proof claims; our search up to the date
of this note does not certify priority for every strengthening.

\begin{thebibliography}{9}
\bibitem{E49} P. Erd\H os, \href{https://www.renyi.hu/~p_erdos/1949-05.pdf}{On some applications of Brun's method}, Acta Sci. Math. (Szeged) \textbf{13} (1949), 57--63.
\bibitem{E65} P. Erd\H os, \href{https://www.renyi.hu/~p_erdos/1965-17.pdf}{Some recent advances and current problems in number theory}, Lectures on Modern Mathematics III (1965), 196--244, esp. p.~209.
\bibitem{Han} K. Han, \emph{A complete candidate proof of Erd\H os Problem 971}, July 2026; \href{https://www.erdosproblems.com/971}{problem 971, proof claim 135}.
\bibitem{Li} S. Li, \href{https://github.com/daizisheng/erdos-971-lean}{A Lean development for Erd\H os 971} (September 2026), statement and blueprint.
\bibitem{Leung} S.-K. Leung, \href{https://doi.org/10.1515/forum-2024-0292}{Moments of primes in progressions to a large modulus}, Forum Math. \textbf{38} (2026), 555--574.
\bibitem{GP} A. Granville and C. Pomerance,
\href{https://math.dartmouth.edu/~carlp/least.pdf}{On the least prime in certain arithmetic progressions}, J. London Math. Soc. (2) \textbf{41} (1990), 193--200.
\bibitem{LiPr} J. Li, K. Pratt, G. Shakan,
\href{https://arxiv.org/abs/1607.02543}{A lower bound for the least prime in an arithmetic progression}, Q. J. Math. \textbf{68} (2017), 729--758.
\bibitem{OAI} OpenAI,
\href{https://github.com/openai/math/tree/main/preprints/Positive-lower-density-of-large-prime-gaps-September-25-2026}{Positive lower density of large prime gaps}, 25 September 2026, esp. Props.~3.1, 4.2, 5.1 and Lemmas 4.1, 5.2--5.4.
\bibitem{BV} E. Bombieri, On the large sieve, Mathematika \textbf{12} (1965), 201--225; see also the maximal-endpoint formulation and log-weight reduction in \cite[Lemma~5.2]{OAI}.
\end{thebibliography}
\end{document}
